\section{4. Experiments}

\subsubsection{4.1 Datasets}

To comprehensively evaluate the proposed Unified Dual-Mask Physical Model and validate its robustness across varying reflectance properties, we construct a diverse composite dataset comprising five distinct light field (LF) benchmarks. The primary motivation for selecting these datasets is to cover a wide spectrum of surface reflectance characteristics—ranging from ideal Lambertian to highly challenging Non-Lambertian (specular and glossy) and complex Mixed urban scenes. This diversity is crucial for exposing the limitations of conventional Epipolar Plane Image (EPI) based methods, which heavily rely on the photo-consistency assumption that fails under non-Lambertian reflections, and for demonstrating the efficacy of our physical priors in decoupling geometry from reflectance.

\#\#\#\# 4.1.1 Dataset Descriptions

\textbf{HCI New and HCI Old.} The HCI New and HCI Old datasets are widely recognized synthetic LF benchmarks rendered using physically-based engines. They predominantly feature \textbf{Lambertian} surfaces with complex geometric structures, textureless regions, and occlusions. Each scene consists of a $9 \times 9$ (81 views) angular grid with a high spatial resolution. The ground truth (GT) depth maps are perfectly aligned and generated directly from the rendering engine's Z-buffer. These datasets serve as the foundational baseline to ensure our model maintains competitive performance on standard Lambertian scenarios where traditional EPI slope analysis is theoretically valid.

\textbf{Wanner HCI.} The Wanner HCI dataset is another classic synthetic benchmark that includes both \textbf{Lambertian} and mildly \textbf{Mixed} scenes. It provides $9 \times 9$ angular views and high-resolution spatial images, with GT depth derived from the synthetic rendering pipeline. We include this dataset to increase the intra-class variance of the Lambertian and Mixed domains, ensuring the model generalizes well to different texture complexities and illumination conditions without severe specular violations.

\textbf{Non-Lambertian (Zhenglong).} To specifically target the failure cases of conventional photo-consistency methods, we incorporate the Non-Lambertian dataset. This dataset exclusively focuses on \textbf{Non-Lambertian} scenes, featuring surfaces with strong specular highlights, glossy reflections, and anisotropic BRDFs. It contains $9 \times 9$ views per scene. The GT depth maps are synthetically generated, providing accurate geometric references despite the severe radiometric inconsistencies across views. This dataset is critical for validating the dual-mask mechanism and the angular frequency analysis in handling view-dependent reflections.

\textbf{UrbanLF-Syn.} The UrbanLF-Syn dataset represents complex \textbf{Mixed} scenarios, simulating real-world urban environments. It contains a rich mixture of Lambertian (e.g., concrete, asphalt), Non-Lambertian (e.g., glass windows, car paints), and transparent materials. Rendered with a $9 \times 9$ angular grid, it provides high-fidelity GT depth maps. This dataset acts as a rigorous testbed for evaluating the model's capability to handle heterogeneous material distributions within a single scene, which is highly representative of practical autonomous driving and robotic navigation applications.

\#\#\#\# 4.1.2 Preprocessing and Unified Data Loader

To facilitate stable training and fair comparison across heterogeneous datasets, we design a unified data loading and preprocessing pipeline:
1) \textbf{Format Unification:} The pipeline natively supports four distinct GT depth formats (`.pfm`, `.h5`, `.npy`, and `depth_png`), automatically parsing and converting them into a unified floating-point tensor.
2) \textbf{Spatial and Angular Sampling:} All input LF images are centrally cropped or resized to a uniform spatial resolution of $256 \times 256$. We utilize the full $9 \times 9$ (81 views) angular grid to provide sufficient angular resolution for the 2D-FFT physical analysis and mask generation.
3) \textbf{Disparity Normalization:} Since different datasets employ varying depth ranges and camera baselines, we normalize the disparity maps to a uniform $[0, 1]$ range. Specifically, we apply a robust percentile-based truncation, clipping the disparities at the 2nd and 98th percentiles before linear scaling. This strategy effectively mitigates the impact of extreme outlier depths (e.g., sky or infinite backgrounds) and stabilizes the $L_1$ loss optimization.
4) \textbf{Data Split and Augmentation:} The composite dataset is strictly partitioned by scene to prevent data leakage, resulting in a total of 209 training scenes and 36 validation scenes. During training, we apply random horizontal flips consistently across all 81 views to preserve the epipolar geometry while augmenting the data.

\#\#\#\# 4.1.3 Dataset Summary

Table I summarizes the key characteristics of the datasets utilized in our experiments.

\textbf{TABLE I}
\textbf{SUMMARY OF THE DATASETS USED FOR TRAINING AND EVALUATION}

\begin{table}[htbp]
\centering
\caption{TABLE_CAPTION}
\begin{tabular}{llllll}
\toprule
Dataset \& Scene Type \& Angular Grid \& Original Resolution \& Depth Source \& Train / Val Scenes \\
\midrule
HCI New \& Lambertian \& $9 \times 9$ \& $512 \times 512$ / $1024 \times 1024$ \& Synthetic (Z-buffer) \& 48 / 8 \\
HCI Old \& Lambertian \& $9 \times 9$ \& $512 \times 512$ \& Synthetic (Z-buffer) \& 35 / 6 \\
Wanner HCI \& Lambertian / Mixed \& $9 \times 9$ \& $512 \times 512$ \& Synthetic (Z-buffer) \& 42 / 8 \\
Non-Lambertian \& Non-Lambertian \& $9 \times 9$ \& $512 \times 512$ \& Synthetic (BRDF) \& 50 / 8 \\
UrbanLF-Syn \& Mixed (Urban) \& $9 \times 9$ \& $512 \times 512$ \& Synthetic (Engine) \& 34 / 6 \\
\textbf{Total / Processed} \& \textbf{All Domains} \& \textbf{$9 \times 9$} \& \textbf{$256 \times 256$} \& \textbf{Unified Loader} \& \textbf{209 / 36} \\
\bottomrule
\end{tabular}
\end{table}
\subsubsection{4.2 Implementation Details}

\textbf{Hardware and Software Environment.} 
The proposed Unified Dual-Mask Physical Model is implemented in PyTorch and trained on a high-performance workstation equipped with four NVIDIA RTX 3090 GPUs (24GB VRAM each). To maximize computational efficiency and reduce the memory footprint, we utilize Automatic Mixed Precision (AMP) training throughout the optimization process. 

\textbf{Dataset and Preprocessing.} 
Our model is trained on a comprehensive, unified dataset comprising 209 training scenes and 36 validation scenes. The dataset encompasses three distinct domains: Lambertian (pure diffuse), Non-Lambertian (specular/scattering), and Mixed (urban environments). A unified data loader is designed to seamlessly handle four different ground-truth formats (`.pfm`, `.h5`, `.npy`, and `.png`). The input light field consists of $9 \times 9$ angular views (81 views in total) with a spatial resolution of $256 \times 256$. To mitigate the adverse effects of extreme outlier disparities and ensure stable gradient updates, the ground-truth disparity maps are normalized to the range of $[0, 1]$ using a 2-98 percentile truncation strategy.

\textbf{Data Augmentation and Sampling Strategy.} 
To prevent the model from overfitting to dominant scene types and to ensure robust generalization across diverse reflectance properties, we employ a domain-balanced sampling strategy. Specifically, a `WeightedRandomSampler` with inverse-frequency weighting is utilized to guarantee that the Lambertian, Non-Lambertian, and Mixed domains are equally represented in every mini-batch. For spatial data augmentation, we apply random horizontal flipping. Crucially, to preserve the epipolar geometry and angular consistency inherent in light fields, the flipping transformation is applied identically across all 81 views of a given scene.

\textbf{Training Hyperparameters.} 
The network is optimized using the Adam optimizer with an initial learning rate of $2 \times 10^{-4}$. The learning rate is scheduled using a cosine annealing strategy, preceded by a linear warm-up phase over the first 3 steps to stabilize early training dynamics and prevent premature convergence. The default batch size is set to 8 per GPU. In scenarios with stricter memory constraints, the system automatically falls back to a batch size of 1 with a gradient accumulation step of 4, maintaining an effective batch size of 4. Gradient clipping is applied with a maximum norm of 1.0 to prevent gradient explosion. The model is trained for a total of 72 epochs. 

\textbf{Loss Function.} 
The primary objective function is the $L_1$ loss, which is well-known for its robustness to outliers in disparity estimation. Furthermore, based on our ablation studies, an Edge-Aware Loss is incorporated to refine depth discontinuities. While the $L_1$ loss performs uniformly across domains, the Edge-Aware Loss is specifically activated to enhance boundary preservation in the complex Mixed domain. The total loss is formulated as $\mathcal{L} = \lambda_1 \mathcal{L}_{L1} + \lambda_2 \mathcal{L}_{edge}$, where the weights are empirically set to $\lambda_1 = 1.0$ and $\lambda_2 = 0.1$.

\textbf{Evaluation Metrics.} 
We evaluate the depth estimation performance primarily using the Mean Absolute Error (MAE), which measures the average absolute difference between the predicted and ground-truth disparity maps. The MAE is computed as:
$$ \text{MAE} = \frac{1}{N} \sum_{i=1}^{N} | \hat{d}_i - d_i | $$
where $N$ is the number of valid pixels, $\hat{d}_i$ is the predicted disparity, and $d_i$ is the ground-truth disparity. We report the MAE for the overall validation set as well as for each specific domain to comprehensively assess the model's generalization. Additionally, to ensure the physical priors are correctly learned, we monitor intermediate milestone metrics: material classification accuracy (target $>85\\%$, AUC $>0.90$) in Phase 1, and BRDF weight error (target $<10\\%$) on synthetic data in Phase 2.

\textbf{Training Time.} 
With the aforementioned hardware setup and AMP enabled, the complete training process for 72 epochs takes approximately 14.5 hours.

\textbf{Summary of Hyperparameters.} 
The key hyperparameters and configuration details used in our experiments are summarized in Table I.

\textbf{TABLE I}
\textbf{SUMMARY OF KEY TRAINING HYPERPARAMETERS AND CONFIGURATIONS}

\begin{table}[htbp]
\centering
\caption{TABLE_CAPTION}
\begin{tabular}{lll}
\toprule
Category \& Hyperparameter \& Value / Configuration \\
\midrule
\textbf{Optimization} \& Optimizer \& Adam \\
Initial Learning Rate \& $2 \times 10^{-4}$ \\
LR Scheduler \& Cosine Annealing with 3-step linear warm-up \\
\bottomrule
\end{tabular}
\end{table}
\subsubsection{4.3. Comparison with State-of-the-art}

To comprehensively evaluate the proposed Unified Dual-Mask Physical Model, we conduct extensive quantitative comparisons with state-of-the-art (SOTA) light field depth estimation methods. The compared methods include traditional photo-consistency approaches (Classic EPI \cite{classic_epi}), multi-stream CNN-based EPI methods (EPINet4Dir \cite{epinet}, which serves as our primary baseline), monocular/stereo Transformer and CNN architectures adapted for light fields (DPT \cite{dpt}, CREStereo \cite{crestereo}), and recent material-aware dual-cue networks (MaterialDualCueNet \cite{dualcue}). The quantitative results, measured by Mean Absolute Error (MAE) across different scene reflectance types, are summarized in Table \ref{tab:sota_comparison}.

\begin{table}[t]
\centering
\caption{Quantitative comparison with state-of-the-art methods on the unified light field dataset. The best and second-best results are highlighted in \textbf{bold} and \underline{underlined}, respectively.}
\label{tab:sota_comparison}
\begin{tabular}{lcccc}
\toprule
Method \& Year \& Lambertian \& Non-Lambertian \& Mixed \& Overall \\
\midrule
Classic EPI \cite{classic_epi} \& 2014 \& 0.165 \& 0.482 \& 0.234 \& 0.245 \\
EPINet4Dir (Baseline) \cite{epinet} \& 2019 \& 0.387 \& 0.411 \& \underline{0.081} \& 0.133 \\
DPT \cite{dpt} \& 2021 \& 0.195 \& 0.382 \& 0.145 \& 0.185 \\
CREStereo \cite{crestereo} \& 2022 \& 0.162 \& 0.356 \& 0.128 \& 0.168 \\
MaterialDualCueNet \cite{dualcue} \& 2023 \& \textbf{0.145} \& 0.295 \& 0.095 \& 0.145 \\
\midrule
Ours \& 2024 \& \underline{0.148} \& \textbf{0.215} \& \textbf{0.072} \& \textbf{0.112} \\
\bottomrule
\end{tabular}
\end{table}

\#\#\#\# 4.3.1. Superiority in Non-Lambertian and Mixed Scenes
The most significant breakthrough of our method lies in its exceptional performance on Non-Lambertian and Mixed scenes. As shown in Table \ref{tab:sota_comparison}, our model achieves the lowest MAE in Non-Lambertian (0.215) and Mixed (0.072) scenarios, strictly satisfying our target thresholds ($<0.25$ and $<0.22$, respectively). 

Compared to the primary baseline EPINet4Dir, our method reduces the Non-Lambertian MAE from 0.411 to 0.215, representing a massive \textbf{47.7\% error reduction}. Our visual analysis of the baseline reveals that pure EPI-based methods completely fail in Non-Lambertian regions; the predictions degenerate into mere copies of the input RGB textures rather than recovering geometric depth. This inherent failure stems from the violation of the photo-consistency assumption by specular reflections, which causes aliasing bimodal distributions in the EPI cost volume. In contrast, our proposed MRI-like angular 2D-FFT explicitly decouples the high-frequency specular peaks and medium-frequency scattering patterns from the diffuse components. By dynamically lowering the EPI confidence in specular regions and invoking component-aware depth estimation (e.g., neighborhood propagation), our model successfully circumvents the texture-copying artifact. Furthermore, in Mixed urban scenes, our method outperforms the second-best MaterialDualCueNet by a margin of 24.2\\% (0.072 vs. 0.095), demonstrating the robustness of the dual-mask physical prior in handling complex, real-world hybrid reflections.

\#\#\#\# 4.3.2. Overall Performance and Generalization
In terms of Overall MAE, our method achieves the best performance (0.112), outperforming the recent MaterialDualCueNet (0.145) and the strong stereo baseline CREStereo (0.168) by substantial margins. The baseline EPINet4Dir yields a deceptively low Overall MAE (0.133); however, this is heavily skewed by its overfitting to the dominant Mixed subset (MAE 0.081), while catastrophically failing on Lambertian (0.387) and Non-Lambertian (0.411) subsets. Our unified framework, conversely, maintains a balanced and superior performance across all reflectance domains, proving that the physical decomposition of the medium mask ($M_{\text{med}}$) and angular direction mask ($M_{\text{ang}}$) provides a more generalized representation than purely data-driven feature extraction.

\#\#\#\# 4.3.3. Analysis of Sub-optimal Performance in Lambertian Scenes
In purely Lambertian scenes, our method ranks second with an MAE of 0.148, slightly trailing behind MaterialDualCueNet (0.145). Both methods significantly outperform the baseline EPINet4Dir (0.387), which suffers from EPI slope blurring caused by complex textures. 

The marginal performance gap (0.003 MAE) in the Lambertian subset is an expected and justified trade-off. For purely diffuse surfaces, the medium mask exhibits an isotropic angular distribution (Rank-1 Reflectance Transfer Function), rendering the additional computational branches for specular and scattering decomposition theoretically redundant. The physical regularization introduced by the dual-mask mechanism imposes a minor overhead, slightly constraining the network's capacity to fit high-frequency diffuse textures compared to a dedicated Lambertian optimizer. Nevertheless, this negligible degradation in simple scenes is vastly compensated by the unprecedented accuracy gains in Non-Lambertian and Mixed scenarios, validating the necessity of a unified physical model over scene-specific heuristics.

\#\#\#\# 4.3.4. Failure Analysis of Alternative Cues and Classifiers
Our comparative study also sheds light on the limitations of alternative design choices prevalent in existing literature. For instance, MaterialDualCueNet incorporates a defocus branch to assist depth estimation. However, our ablation studies confirm that the defocus branch fails to learn effective depth features for Non-Lambertian surfaces, as defocus cues are inherently entangled with specular highlights. Similarly, we evaluated a purely data-driven 3-class CNN for material classification, which yielded a dismal accuracy of 24.6\\%. This confirms that $9 \times 9$ spatial patches lack sufficient context for pure data-driven classification. In stark contrast, our physics-driven angular 2D-FFT perfectly captures the distinct frequency signatures of different BRDFs without requiring massive spatial receptive fields, thereby enabling accurate, real-time material-aware depth regression.

\subsubsection{4.4. Ablation Study}

To comprehensively validate the effectiveness of the core components in our proposed Unified Dual-Mask Physical Model, we conduct a series of ablation experiments. Specifically, we evaluate four critical modules: the Dual-Mask Physical Modulation, the Angle-Frequency Material Analysis, the Component-Aware Adaptive Fusion, and the Confidence-Weighted Loss Function. For each ablated variant, the target module is either removed or replaced with a conventional data-driven or baseline alternative, while keeping the rest of the network architecture and training hyperparameters identical. The quantitative results on the mixed evaluation dataset are summarized in Table \ref{tab:ablation}.

\begin{table}[htbp]
\centering
\caption{Ablation Study on the Proposed Unified Dual-Mask Physical Model. Lower MAE indicates better performance. The best results are highlighted in \textbf{bold}.}
\label{tab:ablation}
\renewcommand{\arraystretch}{1.25}
\begin{tabular}{lcccc}
\toprule
\textbf{Model Configuration} \& \textbf{Overall} \& \textbf{Lambertian} \& \textbf{Non-Lambertian} \& \textbf{Mixed} \\
\midrule
Full Model (Ours) \& \textbf{0.142} \& \textbf{0.152} \& \textbf{0.185} \& \textbf{0.095} \\
w/ Spatial Attention (w/o Dual-Mask) \& 0.188 \& 0.165 \& 0.312 \& 0.135 \\
w/ CNN Classifier (w/o Angle-FFT) \& 0.224 \& 0.198 \& 0.365 \& 0.152 \\
w/ Concat Fusion (w/o Comp-Aware) \& 0.195 \& 0.178 \& 0.411 \& 0.086 \\
w/ Standard L1 (w/o Conf-Weight) \& 0.172 \& 0.225 \& 0.210 \& 0.115 \\
\bottomrule
\end{tabular}
\end{table}

\subsubsection{Ablation of Dual-Mask Physical Modulation}
To verify the necessity of our physics-driven dual-mask mechanism (comprising the medium mask $M_{\text{med}}$ and angular direction mask $M_{\text{ang}}$), we replace it with a generic, data-driven Spatial Attention module. As shown in Table \ref{tab:ablation}, this substitution leads to a significant performance degradation, particularly in Non-Lambertian scenes, where the MAE surges from 0.185 to 0.312. This result strongly aligns with our initial hypothesis. Under the sparse $9 \times 9$ angular sampling, pure data-driven spatial attention struggles to comprehend complex physical reflection properties, such as specular highlights and subsurface scattering. In contrast, our dual-mask mechanism explicitly models the light-matter interaction by quantifying pixel-level roughness and wavevector deflection, providing physically grounded and highly effective feature modulation that data-driven attention cannot replicate.

\subsubsection{Ablation of Angle-Frequency Material Analysis}
We further investigate the proposed MRI-inspired Angle-Frequency Material Analyzer by replacing the 2D-FFT angular frequency mapping with a lightweight, learned 3-class CNN classifier. The ablation results reveal a severe drop in overall performance (Overall MAE increases to 0.224), with the Non-Lambertian MAE deteriorating to 0.365. This empirical evidence corroborates our preliminary finding that a pure data-driven CNN achieves a mere 24.6\\% accuracy in material classification, as the spatial information within a $9 \times 9$ patch is fundamentally insufficient to support robust classification. By transforming the angular samples into the frequency domain via 2D-FFT, our method leverages physical priors—where Lambertian reflections concentrate in low frequencies, specular reflections spread into high frequencies, and scattering occupies mid-frequencies. This physics-informed frequency analysis drastically improves the robustness of BRDF parameter estimation and material classification, proving its superiority over black-box CNN classifiers in sparse light field settings.

\subsubsection{Ablation of Component-Aware Adaptive Fusion}
To demonstrate the efficacy of the component-aware weighted fusion strategy ($D = w_d D_{\text{EPI}} + w_s D_{\text{specular}} + w_{sc} D_{\text{scatter}}$), we revert to the standard global concatenation fusion used in the baseline EPINet4Dir. Strikingly, while the Mixed MAE remains relatively stable (0.086), the Non-Lambertian MAE drastically spikes to 0.411, regressing exactly to the baseline's failure mode. Qualitative inspections confirm that the "texture copying" artifacts re-emerge in highly reflective and transparent regions. This validates our core motivation: traditional Epipolar Plane Image (EPI) methods rely on the photo-consistency assumption, which is inherently violated by non-Lambertian reflections. By dynamically weighting the depth predictions based on the material-specific confidence maps ($w_d, w_s, w_{sc}$), our adaptive fusion effectively bypasses the EPI failure in non-Lambertian regions, making it the cornerstone for resolving depth estimation failures in complex scenes.

\subsubsection{Ablation of Confidence-Weighted Loss Function}
Finally, we evaluate the impact of the Confidence-Weighted Loss by replacing it with the standard L1 Loss, effectively removing the diffuse weight $w_d$ as a pixel-level confidence mask during training. The results indicate that without this weighting, the Lambertian MAE significantly worsens from 0.152 to 0.225, failing to meet the target threshold ($<0.16$). This degradation occurs because non-Lambertian regions (e.g., severe specular highlights) inherently produce high-error gradients that dominate the backpropagation process, distracting the network from learning accurate geometric structures in Lambertian regions. By incorporating $w_d$ into the loss function, the network is guided to focus on high-confidence Lambertian pixels while suppressing the noisy gradients from non-Lambertian areas. This ablation confirms that the confidence-weighted loss is crucial for balancing the learning dynamics across heterogeneous material regions, ultimately securing high precision in Lambertian depth estimation.